% !PRINTED 269
\weberpage{269}
It should still be mentioned that the discussion carried out in §\,69, 4., 5. for the invariant equation in the case of a composite degree of transformation holds unchanged for the \gesperrt{Schlaefli} modular equations as well, from which we may conclude that all these equations have rational numerical coefficients.

\begin{center}
\textbf{§\,75.\quad \gesperrt{The Irrational Forms of the Modular Equations.}}
\end{center}

By applying the same procedure, far simpler forms can be given to the transformation equations. The forms of equation with which we shall now occupy ourselves contain the three functions $f$, $f_1$, $f_2$ at once; but since, by §\,34, (11), (12), two of these functions can be expressed by the third, two of them can be eliminated, and in this way one can arrive at the equations of the preceding section. If one substitutes for $f_1$, $f_2$ the expressions in terms of $f$, or for the three functions $f$, $f_1$, $f_2$ the expressions in terms of $k^2$, then radical signs occur, and this explains the name of these equations.

We set, as above,
\weq{1}{\begin{aligned}
u &= f(\omega), & v &= f\Bigl(\frac{c+\partial\,\omega}{a}\Bigr),\\
u_1 &= f_1(\omega), & v_1 &= \Bigl(\frac{2}{a}\Bigr) f_1\Bigl(\frac{c+\partial\,\omega}{a}\Bigr),\\
u_2 &= f_2(\omega), & v_2 &= \Bigl(\frac{2}{\partial}\Bigr) f_2\Bigl(\frac{c+\partial\,\omega}{a}\Bigr),\\
u\,u_1\,u_2 &= \sqrt{2}, & v\,v_1\,v_2 &= \Bigl(\frac{2}{n}\Bigr)\sqrt{2},
\end{aligned}}
and obtain the following corresponding permutations:
\weq{2}{\begin{array}{llll}
\omega, & uv, & u_1v_1, & u_2v_2,\\
-\dfrac{1}{\omega}, & \varrho\,uv, & \varrho\,u_2v_2, & \varrho\,u_1v_1,\\
\omega+1, & e^{-\frac{(n+1)\pi i}{24}}\sigma\,u_1v_1, & e^{-\frac{(n+1)\pi i}{24}}\sigma\,uv, & e^{\frac{(n+1)\pi i}{12}}\sigma\,u_2v_2,
\end{array}}
% Source scan and the row pattern confirm the exponent numerator (n+1).
where $\varrho$, $\sigma$ are the third roots of unity defined above [§\,73, (4)].

We distinguish three different cases according to the behaviour of $n$ with respect to the modulus 8.

1. $n+1\equiv 0$, (mod 8).
